%===========================================================
\chapter{Predictive Modelling, Constrained Selection and Biological Interpretation}
\label{chap:model_training_evaluation}
%===========================================================

%-----------------------------------------------------------
\section{Introduction}
%-----------------------------------------------------------

Chapter~\ref{chap:feature_selection} established a statistically stable and 
structurally diversified CpG panel of fixed cardinality ($K = 5{,}000$) for 
each dataset. The proposed pipeline integrated empirical variability screening, 
biologically weighted stability selection, correlation-based redundancy pruning, 
and genomic diversification, ensuring that the resulting signatures are 
reproducible within each cohort and structurally constrained by design. 
The inter-dataset evaluation of Section~5.4 further revealed that cross-cohort 
replicability is partial and asymmetric: GSE69914 and GSE225845 share a 
concordant drift orientation and support bidirectional zero-shot transfer, 
whereas GSE287331 exhibits a systematically inverted drift axis, yielding 
sub-random transfer performance against both partner cohorts.

These structural properties define the design space for the analyses developed 
in the present chapter. The selected CpG panels are characterised by statistical 
robustness and biological diversification, but their capacity to discriminate 
Normal from Adjacent tissue under supervised modelling has not yet been 
assessed. Leakage-controlled feature selection guarantees that the chosen loci 
are statistically stable and non-redundant; it does not, by itself, imply that 
they encode sufficient discriminative information for a classifier to operate 
reliably on held-out data.

The objectives of this chapter are threefold. First, intra-dataset 
classification performance is quantified to verify whether the selected CpG 
panels support accurate discrimination of Normal and Adjacent tissue within 
each cohort. A linear Support Vector Machine is adopted as the primary 
classifier for GSE225845 and GSE287331, a $k$-Nearest Neighbours introduced for 
GSE69914 to assess the degree of linear separability. Second, cross-dataset 
transferability is examined by training on one cohort and evaluating directly 
on another, with no information from the target dataset entering the training 
procedure. This zero-leakage design quantifies the extent to which the inferred 
epigenetic drift axis generalises across independent cohorts, and whether 
performance degradation is systematically associated with the directional 
concordance structure documented in Chapter~\ref{chap:feature_selection}. 
Third, a constrained combinatorial compression step is introduced to derive a 
compact 50-CpG subset through a mixed-integer linear programming formulation 
inspired by the classical 0--1 knapsack paradigm. This step investigates 
whether predictive structure and biological diversity can be jointly preserved 
under severe dimensional reduction, subject to explicit constraints on 
cardinality, chromosomal distribution, correlation structure, and 
methylation directionality.

Finally, the biological coherence of the compressed panels is evaluated through 
CpG-to-gene mapping, cross-referenced against the COSMIC Cancer Gene Census (CGC, version 103, GRCh38)~\cite{sondka2018cosmic,cosmic_cgc_v103}, 
and pathway-level enrichment analysis. This interpretive step closes the loop 
between statistical discrimination and functional relevance, assessing whether 
the loci selected by a purely data-driven optimisation converge on genomic 
regions with established roles in cancer-related epigenetic deregulation.

Together, these analyses progress from statistically robust feature selection 
to supervised evaluation, external validation, and structured biological 
interpretation. The limitations that emerge — in particular the failure of 
linear $\Delta\beta$ signatures to generalise across cohorts with divergent 
drift orientations — directly motivate the representation-learning approach 
developed in Chapter~\ref{chap:CpGPT}.


%-----------------------------------------------------------
\section{Predictive Framework}
%-----------------------------------------------------------

Following feature selection, each dataset is restricted to the 
corresponding panel of $K = 5{,}000$ CpGs identified in 
Chapter~\ref{chap:feature_selection}. 
Supervised baseline classifiers are applied independently 
within each cohort to assess whether the structurally stable 
and biologically diversified signatures encode sufficient 
discriminative information to separate Normal and Adjacent tissue.

%-----------------------------------------------------------
\subsection{Intra-Dataset Supervised Modelling}
%-----------------------------------------------------------

\paragraph{Data partitioning and leakage control}
For each dataset, samples are partitioned into stratified 
training and test sets (80/20 split), preserving the 
Normal--Adjacent class proportions. Standardisation is fitted 
exclusively on the training set and then applied to the held-out 
test set without re-estimation. Feature selection is performed 
strictly within training folds in Chapter~\ref{chap:feature_selection}, 
ensuring complete leakage control.

\paragraph{Linear Support Vector Machine}
A linear Support Vector Machine (SVM) is adopted as the primary 
classifier~\cite{cortes1995svm}. Given training data 
$\{(x_i, y_i)\}_{i=1}^n$ with $x_i \in \mathbb{R}^p$ and 
$y_i \in \{-1,+1\}$, the primal formulation is

\[
\min_{w,b,\xi} \quad 
\frac{1}{2}\lVert w\rVert^2 + C \sum_{i=1}^n \xi_i
\quad
\text{s.t.}\quad
y_i(w^\top x_i + b) \ge 1 - \xi_i,\ \xi_i \ge 0.
\]

The decision function is $f(x)=w^\top x + b$, with class assignment 
given by $\mathrm{sign}(f(x))$. In this work, the model is trained on 
standardised $M$-values with fixed hyperparameters ($C=0.5$, 
$\texttt{max\_iter}=10{,}000$) to prioritise robustness over tuning.

\paragraph{$k$-Nearest Neighbours (GSE69914 only)}
For GSE69914, a $k$-Nearest Neighbours (kNN) classifier is also evaluated 
as a local non-parametric baseline~\cite{cover1967nn}. For a query $x$, 
let $\mathcal{N}_k(x)$ denote the set of the $k$ closest training samples 
under Euclidean distance. The predicted label is

\[
\hat{y}(x)=\operatorname{mode}\{y_i:\ x_i \in \mathcal{N}_k(x)\}.
\]

The method is applied on standardised features with fixed $k=5$.

\paragraph{Performance metrics}
AUC is computed as defined in Chapter~\ref{chap:feature_selection}. 
Balanced accuracy, MCC, precision, recall, and F1-score are reported to 
complement AUC under potential class imbalance. Balanced accuracy is

\[
\mathrm{BAcc}=\frac{1}{2}\left(
\frac{\mathrm{TP}}{\mathrm{TP}+\mathrm{FN}}+
\frac{\mathrm{TN}}{\mathrm{TN}+\mathrm{FP}}
\right),
\]

while MCC is

\[
\mathrm{MCC}=
\frac{\mathrm{TP}\cdot\mathrm{TN}-\mathrm{FP}\cdot\mathrm{FN}}
{\sqrt{(\mathrm{TP}+\mathrm{FP})(\mathrm{TP}+\mathrm{FN})
(\mathrm{TN}+\mathrm{FP})(\mathrm{TN}+\mathrm{FN})}}.
\]

Finally,
\[
\mathrm{Precision}=\frac{\mathrm{TP}}{\mathrm{TP}+\mathrm{FP}},
\qquad
\mathrm{Recall}=\frac{\mathrm{TP}}{\mathrm{TP}+\mathrm{FN}},
\qquad
\mathrm{F1}=\frac{2\,\mathrm{Precision}\,\mathrm{Recall}}
{\mathrm{Precision}+\mathrm{Recall}}.
\]

%-----------------------------------------------------------
\subsection{Cross-Dataset Evaluation Protocol}
%-----------------------------------------------------------

This subsection formalises the cross-cohort transferability 
framework adopted to assess whether the inferred epigenetic 
drift axis generalises across independent datasets.

\paragraph{Feature space alignment}

Cross-dataset evaluation requires strict alignment of the 
feature space between source and target cohorts. 
For each ordered pair of datasets $(A,B)$, the CpG panel is 
restricted to the intersection of the effective 5,000-CpG sets 
used during intra-dataset training in cohort $A$. 
Only CpGs present in both datasets and retained in the 
final feature selection of the source cohort are considered.

Column ordering is explicitly harmonised to ensure identical 
feature indexing across matrices. No re-selection or 
feature re-ranking is performed on the target cohort. 
All transformations (including the $\beta \to M$ mapping, 
when required) are applied consistently across datasets 
prior to modelling.

\paragraph{Train-on-A / Test-on-B design}

External validation is implemented under a strict 
train-on-$A$ / test-on-$B$ paradigm. 
The classifier is trained exclusively on the training split 
of cohort $A$, using the aligned CpG space described above. 
No samples from cohort $B$ enter the fitting procedure at any stage.

Feature standardisation is fitted on the source training set 
and subsequently applied, without re-estimation, to the 
target dataset. A linear Support Vector Machine is trained 
on the source cohort and probability calibration is performed 
via sigmoid calibration (three-fold internal cross-validation) 
within the source training data only. 
The calibrated decision function is then evaluated directly 
on cohort $B$.

No hyperparameter tuning, retraining, or adaptation 
is performed on the target cohort. 
This design enforces complete separation between 
source and target distributions and quantifies 
true zero-shot transferability.

\paragraph{Evaluation criteria}

Discriminative performance is quantified using the same 
metrics adopted for intra-dataset validation, including 
AUC, balanced accuracy, and Matthews correlation coefficient. 
Performance is reported both on the internal test split 
of cohort $A$ and on the external cohort $B$.

Under distributional shift, degradation in AUC or balanced 
accuracy is interpreted as evidence of limited transferability 
of the learned separating hyperplane. 
In particular, systematic performance inversion is examined 
in relation to the directional concordance structure 
identified in Chapter~\ref{chap:feature_selection}. 
External performance below random expectation 
(AUC $\approx 0.5$) is considered indicative of 
either weak shared signal or reversed drift orientation 
between cohorts.

%-----------------------------------------------------------
\subsection{Constrained Subset Optimisation: Knapsack Formulation}
%-----------------------------------------------------------

This subsection formalises the compression of the 
5,000-CpG panel into a fixed-cardinality subset of size 
$K = 50$ through a constrained combinatorial optimisation framework. 
The formulation is inspired by the classical $0$--$1$ knapsack 
problem originally discussed in the context of discrete-variable 
extremum problems by Dantzig~\cite{dantzig1957}, and extended here 
to incorporate multiple structural and biological constraints.

While the classical knapsack maximises a linear utility function 
under a single capacity constraint, the present formulation 
adopts a multi-constraint mixed-integer linear programming (MILP) 
structure tailored to genomic panel design. Similar optimisation-based 
approaches have been explored in marker and SNP panel selection 
problems~\cite{wu2016snp,nishiyama2022ilp}, supporting the use 
of binary decision variables and explicit structural constraints 
in genomics.

\paragraph{Objective function}

Let $\mathcal{J}$ denote the set of 5,000 CpGs retained after 
feature selection. For each CpG $j \in \mathcal{J}$, 
a binary decision variable $x_j \in \{0,1\}$ indicates inclusion 
in the compressed panel.

Each CpG is assigned a composite score
\[
v_j = 
w_{\mathrm{stat}} \,\tilde{s}_j 
+ w_{\Delta M} \,\tilde{d}_j 
+ w_{\mathrm{stab}} \,\tilde{r}_j
+ w_{\mathrm{COSMIC}} \,\mathbf{1}_{\{j \in \mathrm{CGC}\}},
\]
where $\tilde{s}_j$, $\tilde{d}_j$, and $\tilde{r}_j$ denote 
normalised statistical ranking components 
(classifier weight magnitude, absolute effect size, and stability score, 
respectively), and $\mathbf{1}_{\{j \in \mathrm{CGC}\}}$ indicates 
whether the gene mapped to CpG $j$ appears in the 
COSMIC Cancer Gene Census~\cite{sondka2018cosmic}.

The optimisation objective is

\[
\max_{x} \quad \sum_{j \in \mathcal{J}} v_j x_j
+ \mu \, \Phi(x),
\]

where $\Phi(x)$ encodes diversification incentives 
and $\mu \ge 0$ controls the trade-off between 
pure statistical ranking and structural balance. 
A sweep over $\mu$ is performed to identify 
a stable trade-off region.

\paragraph{Cardinality and structural constraints}

The optimisation is subject to the following constraints:

\begin{itemize}
\item \textbf{Fixed cardinality:}
\[
\sum_{j \in \mathcal{J}} x_j = K.
\]

\item \textbf{Chromosome caps:}
for each chromosome $c$,
\[
\sum_{j \in \mathcal{J}_c} x_j \le \alpha_{\mathrm{chr}} K,
\]
where $\mathcal{J}_c$ denotes CpGs located on chromosome $c$.

\item \textbf{Genomic context limits:}
upper bounds are imposed on the proportion of CpGs 
belonging to specific contexts (e.g., OpenSea, Island).

\item \textbf{Correlation exclusion:}
for CpG pairs $(i,j)$ with 
$|r_{ij}| \ge 0.85$,
\[
x_i + x_j \le 1,
\]
preventing the simultaneous selection of highly 
correlated loci.

\item \textbf{Gene-level diversity:}
linking variables ensure that no gene is represented 
by more than a fixed number of CpGs, and that 
a minimum number of distinct genes is selected.

\item \textbf{Directional balance:}
a minimum number of hypermethylated loci is enforced 
to prevent polarity collapse.
\end{itemize}

These constraints extend the classical single-capacity 
knapsack formulation~\cite{dantzig1957} into a structured 
multi-constraint MILP suitable for genomic panel compression.

\paragraph{Optimisation strategy}

The resulting problem is solved using a mixed-integer 
linear programming solver (CBC backend via PuLP). 
Given the moderate dimensionality ($|\mathcal{J}| = 5{,}000$), 
exact optimisation is computationally feasible. 

To assess stability of the solution with respect to 
the diversification parameter $\mu$, a parameter sweep 
is conducted and the Pareto trade-off curve between 
statistical score and diversification component 
is analysed. The final subset is selected at the 
empirical knee point of this curve.

The formulation therefore combines classical 
binary combinatorial optimisation~\cite{dantzig1957} 
with recent optimisation-based approaches to marker 
panel design in genomics~\cite{wu2016snp,nishiyama2022ilp}, 
yielding a compact CpG panel that jointly preserves 
predictive relevance and biological structure.

%-----------------------------------------------------------
\subsection{Gene Set and Functional Enrichment Analysis}
%-----------------------------------------------------------

This subsection evaluates the biological coherence of the 
CpG subsets selected through supervised modelling and 
constrained optimisation. The objective is to determine 
whether loci identified purely through statistical and 
structural criteria converge on genes with established 
roles in cancer biology.

\paragraph{CpG-to-gene mapping}

CpG loci are annotated to genes using the Illumina 
EPIC manifest (GRCh38 assembly), relying on the 
\texttt{UCSC\_RefGene\_Name} field. 
When multiple gene symbols are associated with a single CpG, 
all mapped genes are retained. 
Gene symbols are harmonised to official HGNC nomenclature 
prior to downstream analysis.

Let $\mathcal{G}$ denote the set of unique genes mapped 
from the selected CpG subset. 
Gene-level analyses are conducted on $\mathcal{G}$ 
after removal of duplicated mappings.

\paragraph{Enrichment testing}

Functional coherence is assessed using over-representation 
analysis (ORA). 
Let $\mathcal{U}$ denote the background universe of genes 
mapped from the full 5,000-CpG panel of the corresponding cohort. 
For a given reference gene set $\mathcal{S}$, 
the enrichment of $\mathcal{G}$ in $\mathcal{S}$ 
is evaluated via the hypergeometric test:

\[
p =
\sum_{k = k_0}^{|\mathcal{G}|}
\frac{
\binom{|\mathcal{S}|}{k}
\binom{|\mathcal{U}| - |\mathcal{S}|}{|\mathcal{G}| - k}
}{
\binom{|\mathcal{U}|}{|\mathcal{G}|}
},
\]

where $k_0 = |\mathcal{G} \cap \mathcal{S}|$ denotes 
the observed overlap. 

Multiple hypothesis testing correction is performed 
using the Benjamini--Hochberg procedure to control 
the false discovery rate (FDR). 
Adjusted $q$-values below $0.05$ are considered significant.

\paragraph{Reference databases}

External biological knowledge is incorporated through 
curated gene sets. 
In particular, selected genes are cross-referenced 
against the COSMIC Cancer Gene Census (CGC) database 
(version 103, GRCh38 assembly)~\cite{sondka2018cosmic}, 
which catalogues genes with documented roles in oncogenesis.

Additional curated gene lists relevant to field cancerisation 
and early epigenetic deregulation are considered when 
interpreting enrichment results. 
The integration of these external resources enables 
assessment of whether statistically derived CpG panels 
converge on genes implicated in tumour initiation, 
chromatin regulation, and cancer-associated signalling pathways.

%-----------------------------------------------------------
\section{Intra-Dataset Results}
%-----------------------------------------------------------

\subsection{Dataset GSE69914}

\subsection{Dataset GSE225845}

\subsection{Dataset GSE287331}

%-----------------------------------------------------------
\section{Inter-Dataset Results}
%-----------------------------------------------------------

%-----------------------------------------------------------
\subsection{Cross-Cohort Transferability of the 5,000-CpG Panel}
%-----------------------------------------------------------

External validation results across all train/test combinations.

%-----------------------------------------------------------
\subsection{Directional Concordance and Performance Degradation}
%-----------------------------------------------------------

Analysis of performance variation in relation to
inter-dataset $\Delta\beta$ concordance.

%-----------------------------------------------------------
\subsection{Cross-Dataset Behaviour of the 50-CpG Subset}
%-----------------------------------------------------------

Evaluation of compressed signature under distributional shift.

%-----------------------------------------------------------
\section{Synthesis and Implications}
%-----------------------------------------------------------

This section summarises predictive findings,
identifies structural limitations of linear
$\Delta\beta$-based signatures,
and motivates the transition toward
representation-learning approaches
in the subsequent chapter.